% !TeX program = xelatex
\documentclass[12pt,spanish,letterpaper,oneside]{article}
\PassOptionsToPackage{es-noshorthands}{babel}
\def\documenttitle{Resolución de caso ético-jurídico sobre la crisis del agua y aprovechamientos ilícitos en la cuenca citrícola de Nuevo León}
\def\documentsubtitle{Evidencia 1 --- Reporte de caso en contexto ético, jurídico y democrático (Fase 1)}
\def\documentsubject{Ética y Cultura de la Legalidad --- Ingeniero Agrónomo}
\def\documentauthor{Carlos Mauricio Carvajal Coronado}
\def\coursename{Ética y Cultura de la Legalidad}
\def\coursecode{REF\#261013755115}
\def\universityname{Universidad Autónoma de Nuevo León}
\def\universityfaculty{Facultad de Agronomía}
\def\universitydepartment{Ingeniero Agrónomo}
\def\universitydepartmentimage{departamentos/logos-institucionales.pdf}
\def\universitydepartmentimagecfg{height=2.1cm}
\def\universitylocation{General Escobedo / Marín, Nuevo León, México}
\def\authortable{%
	\begin{Form}
	\begin{tabular}{@{}ll@{}}
		\textbf{Estudiante:} & Carlos Mauricio Carvajal Coronado\\[0.35cm]
		\textbf{Matrícula:} & 1977875\\[0.35cm]
		\textbf{Docente:} & Cátedra de Formación General Universitaria\\[0.35cm]
		\textbf{Unidad de Aprendizaje:} & Ética y Cultura de la Legalidad (Evidencia 1 \textbar\ 8 Pts)\\[0.5cm]
		\multicolumn{2}{l}{Fecha de realización: 14 de octubre de 2026}\\
		\multicolumn{2}{l}{\universitylocation}
	\end{tabular}
	\end{Form}
}

\usepackage{fontspec}
\setmainfont{Arial}
\setsansfont{Arial}
\AddToHook{file/config.tex/after}[aulatex-etica]{%
  \def\documentfontsize{12}
  \def\documentinterline{1}
  \def\documentparindent{0}
  \def\documentparskip{5}
  \def\fontdocument{custom}
  \def\fonttypewriter{custom}
  \def\stylecitereferences{custom}
  \def\portraitstyle{style1}
  \def\hfstyle{style1}
  \def\hfwidthwrap{true}
  \def\hfwidthtitle{0.44}
  \def\hfwidthcourse{0.52}
  \def\predocpageromannumber{false}
  \def\pagemarginleft{2.5}
  \def\pagemarginright{2.5}
  \def\pagemargintop{2.5}
  \def\pagemarginbottom{2.5}
  \def\firstpagemargintop{3.8}
  \def\sectionfontsize{\large}
  \def\sectionspacingtop{8}
  \def\sectionspacingbottom{4}
  \def\ssectionfontsize{\normalsize}
  \def\ssectionspacingtop{6}
  \def\ssectionspacingbottom{3}
  \def\defaultimagefolder{UANL/img/}
  \def\usepdfmetadata{false}
}
% Template:     Informe LaTeX
% Documento:    Núcleo del template
% Versión:      DEV
% Codificación: UTF-8
%
% Autor: Pablo Pizarro R.
%        pablo@ppizarror.com
%
% Manual template: [https://latex.ppizarror.com/informe]
% Licencia MIT:    [https://opensource.org/licenses/MIT]

% -----------------------------------------------------------------------------
% CONFIGURACIONES
% -----------------------------------------------------------------------------
% Definiciones previas
\input{src/defs}

% Ajustes usuario
\input{src/config}

% -----------------------------------------------------------------------------
% IMPORTACIÓN DE LIBRERÍAS
% -----------------------------------------------------------------------------
\input{src/env/imports}

% -----------------------------------------------------------------------------
% IMPORTACIÓN DE FUNCIONES Y ENTORNOS
% -----------------------------------------------------------------------------
\input{src/cmd/core}
\input{src/cmd/math}
\input{src/cmd/equation}
\input{src/cmd/image}
\input{src/cmd/title}
\input{src/cmd/other}
\input{src/cmd/column}
\input{src/env/environments}

% -----------------------------------------------------------------------------
% IMPORTACIÓN DE ESTILOS
% -----------------------------------------------------------------------------
\input{src/style/code}
\input{src/style/other}

% -----------------------------------------------------------------------------
% CONFIGURACIÓN INICIAL DEL DOCUMENTO
% -----------------------------------------------------------------------------
\input{src/cfg/init}

% -----------------------------------------------------------------------------
% PORTADA
% -----------------------------------------------------------------------------
\input{src/page/portrait}
\input{src/page/portrait_config}

% -----------------------------------------------------------------------------
% CONFIGURACIÓN DE PÁGINA Y ENCABEZADOS
% -----------------------------------------------------------------------------
\input{src/cfg/page}

% -----------------------------------------------------------------------------
% TABLA DE CONTENIDOS - ÍNDICE
% -----------------------------------------------------------------------------
\input{src/page/index}

% -----------------------------------------------------------------------------
% CONFIGURACIONES FINALES
% -----------------------------------------------------------------------------
\input{src/cfg/final}

\RemoveFromHook{file/config.tex/after}[aulatex-etica]


\usepackage{tabularx}
\usepackage{booktabs}
\usepackage{amsmath,amssymb}
\usepackage[backend=biber,style=apa,sorting=nyt]{biblatex}
\definecolor{uanlblue}{RGB}{0,55,108}
\definecolor{tableblue}{RGB}{233,240,247}
\def\portraittitlecolor{uanlblue}
\def\sectioncolor{uanlblue}
\def\ssectioncolor{uanlblue}
\def\documenttitlehf{Ética y Legalidad: Evidencia 1 --- Resolución de Caso Hídrico}
\hypersetup{colorlinks=true,linkcolor=uanlblue,urlcolor=uanlblue,citecolor=uanlblue,pdfauthor={Carlos Mauricio Carvajal Coronado},pdftitle={Resolución de caso ético-jurídico sobre la crisis del agua en Nuevo León},pdfsubject={Evidencia 1 - Ética y Legalidad UANL}}
\addbibresource{UANL/ingeniero-agronomo/etica-y-cultura-de-la-legalidad/etica-y-cultura-de-la-legalidad.bib}
\newcolumntype{Y}{>{\raggedright\arraybackslash}X}
\setlength{\emergencystretch}{2.5em}

\begin{document}
\hypersetup{pageanchor=false}
\begingroup
\fontsize{14}{17}\selectfont
\renewcommand{\Huge}{\fontsize{14}{17}\selectfont\bfseries}
\renewcommand{\large}{\fontsize{14}{17}\selectfont}
\templatePortrait
\clearpage
\endgroup
\hypersetup{pageanchor=true}
\templatePagecfg
\singlespacing
\setlength{\tabcolsep}{8pt}
\renewcommand{\arraystretch}{1.18}

\section*{Resumen}
En el presente trabajo se formaliza la Evidencia 1 de la Unidad de Aprendizaje de Ética y Cultura de la Legalidad mediante el análisis exhaustivo y la propuesta resolutiva de un caso real de problemática social y ambiental: el desvío y extracción clandestina de aguas superficiales y subterráneas en la cuenca citrícola de Linares y Montemorelos, Nuevo León. A través de la perspectiva teórica de la ética cívica de mínimos de Adela Cortina, la noción de justicia distributiva y los fundamentos del Estado constitucional de derecho, se desglosan las transgresiones al Artículo 4to y 27 de la Constitución Política de los Estados Unidos Mexicanos y a la Ley de Aguas Nacionales. Se formula un dictamen resolutivo que armoniza la sanción punitiva con la recomposición del tejido social agrario y la participación ciudadana en consejos de cuenca.

\vspace{0.4cm}
\noindent\textbf{Palabras clave:} Ética cívica, cultura de la legalidad, derecho al agua, estado de derecho, justicia distributiva, cuenca citrícola.

\section{Presentación del Caso de Estudio}
Durante los ciclos agrícolas recientes, la región citrícola de Nuevo León (comprendida por los municipios de Montemorelos, Linares, Hualahuises y General Terán) experimentó tensiones socioambientales agudas derivadas del estrés hídrico. En inspecciones conjuntas entre usuarios de distritos de riego y autoridades federales, se detectaron más de 80 tomas clandestinas en los cauces de los ríos San Fernando y Ramos, así como represas particulares construidas sin concesión de la Comisión Nacional del Agua (CONAGUA).
Mientras grandes predios comerciales aseguraban volúmenes hídricos para riego suntuario y huertas privadas mediante pozos ilegales, ejidatarios y comunidades ribereñas sufrieron el desecamiento de sus norias de abasto y la pérdida total de cosechas de cítricos.

\section{Análisis Ético, Filosófico y Deontológico}
El dilema ético subyacente enfrenta el utilitarismo egoísta de corto plazo frente a la ética de la responsabilidad y el bien común:
\begin{itemize}
    \item \textbf{Ética Cívica y Mínimos de Justicia:} Siguiendo a \textcite{cortina2013}, una sociedad democrática pluralista se sustenta en una ``ética de mínimos'' irrenunciables: el respeto a la dignidad humana, la equidad en el acceso a los bienes primarios y la no dominación. La apropiación indebida de un bien común como el agua rompe la simetría moral ciudadana.
    \item \textbf{Cultura de la Legalidad:} Como subraya la \textcite{cndh2020}, la cultura de la legalidad no es la mera obediencia pasiva por temor al castigo, sino la convicción interna de que las normas jurídicas protegen los derechos fundamentales de todos. El incumplimiento normalizado del marco hídrico genera un estado de anomia social donde prevalece la ley del más fuerte.
    \item \textbf{Justicia Distributiva y Bien Común:} La distribución de los recursos hídricos debe regirse por criterios de equidad y sustentabilidad intergeneracional, subordinando el interés de lucro privado a la subsistencia colectiva \parencite{giusti2018}.
\end{itemize}

\section{Marco Jurídico y Vulneración del Estado de Derecho}
El caso evidencia la violación de disposiciones del orden constitucional y legal federal:
\begin{enumerate}
    \item \textbf{Artículo 4to Constitucional:} Garantiza el derecho humano de toda persona al acceso, disposición y saneamiento de agua para consumo personal y doméstico en forma suficiente, salubre y asequible \parencite{cpeum2024}.
    \item \textbf{Artículo 27 Constitucional:} Establece que las aguas de los ríos, manantiales y cuencas son propiedad inalienable e imprescriptible de la Nación, sujetas a concesión regulada por el Ejecutivo Federal \parencite{cpeum2024}.
    \item \textbf{Ley Agraria y Ley de Aguas Nacionales:} Sancionan penal y administrativamente la perforación no autorizada de pozos y la obstrucción de escurrimientos en cauces federales \parencite{ley_agraria2022}.
\end{enumerate}

\section{Resolución del Caso y Recomendaciones Estratégicas}
Para dirimir el conflicto desde los principios de la convivencia democrática y el Estado de derecho, se establece el siguiente esquema resolutivo:

\begin{table}[htbp]
\centering
\small
\caption{Matriz resolutiva del conflicto hídrico en la cuenca citrícola.}
\label{tab:resolucion-agua}
\begin{tabularx}{\textwidth}{@{}l Y Y@{}}
\toprule
\textbf{Eje de Acción} & \textbf{Medidas Concretas} & \textbf{Fundamento Ético-Jurídico} \\
\midrule
\textbf{Legal / Coercitivo} & Clausura inmediata de tomas clandestinas, demolición de represas ilegales e inicio de carpetas de investigación por la FGR. & Restitución del Estado de derecho y erradicación de la impunidad \parencite{cndh2020}. \\
\addlinespace
\textbf{Democrático / Social} & Reactivación del Comité Técnico de Aguas Subterráneas (COTAS) con representación paritaria de ejidatarios y pequeños propietarios. & Principio de deliberación democrática y justicia participativa \parencite{cortina2013}. \\
\addlinespace
\textbf{Técnico Agronómico} & Implementación obligatoria de macromedidores telemétricos y transición a riego presurizado por goteo con subsidios públicos focalizados. & Eficiencia biofísica, solidaridad intergeneracional y responsabilidad profesional \parencite{uanl_codigo2021}. \\
\bottomrule
\end{tabularx}
\end{table}

\section{Conclusiones}
1. La resolución ética de los conflictos agrarios exige el apego irrestricto al Estado de derecho, garantizando que el agua sea administrada como un patrimonio ecológico y social de la Nación.
2. Como futuros ingenieros agrónomos formados en la UANL, es deber ético ineludible rechazar cualquier complicidad en obras hidráulicas no autorizadas y promover tecnologías que dignifiquen la vida comunitaria rural.

\newpage
\printbibliography[title={Referencias bibliográficas}]
\end{document}

